\documentclass[12pt]{article}

% ─── Packages ────────────────────────────────────────────────────────────────
\usepackage[utf8]{inputenc}
\usepackage[T1]{fontenc}
\usepackage[spanish]{babel}
\usepackage{geometry}
\usepackage{amsmath, amssymb}
\usepackage{booktabs}
\usepackage{array}
\usepackage{graphicx}
\usepackage{hyperref}
\usepackage{xcolor}
\usepackage{listings}
\usepackage{float}
\usepackage{caption}
\usepackage{subcaption}
\usepackage{multirow}
\usepackage{enumitem}
\usepackage{titlesec}
\usepackage{fancyhdr}

\geometry{margin=2.5cm}

% ─── Colors ──────────────────────────────────────────────────────────────────
\definecolor{codeblue}{RGB}{0,102,204}
\definecolor{codegray}{RGB}{128,128,128}
\definecolor{codered}{RGB}{180,0,0}

% ─── Listings ────────────────────────────────────────────────────────────────
\lstset{
  basicstyle=\ttfamily\small,
  breaklines=true,
  frame=single,
  numbers=left,
  numberstyle=\tiny\color{codegray},
  keywordstyle=\color{codeblue},
  commentstyle=\color{codegray}\itshape,
}

% ─── Header / Footer ─────────────────────────────────────────────────────────
\pagestyle{fancy}
\fancyhf{}
\rhead{GPU Kernel Generation — Reporte Experimental}
\lhead{\leftmark}
\cfoot{\thepage}

% ─── Title ───────────────────────────────────────────────────────────────────
\title{
  \textbf{Generación de Kernels GPU con SLM Aumentado:}\\
  \large Diseño Experimental y Análisis Estadístico\\[0.5em]
  \normalsize Comparación contra Modelos Baseline
}
\author{Cesar Alan Silva Ramos}
\date{\today}

% ─────────────────────────────────────────────────────────────────────────────
\begin{document}
\maketitle
\tableofcontents
\newpage

% =============================================================================
% PARTE I — PRE-REGISTRO
% =============================================================================
\part{Pre-registro Experimental}
\textit{Esta sección fue definida antes de ejecutar los experimentos comparativos.}

% -----------------------------------------------------------------------------
\section{Objetivo Experimental}

El objetivo de este experimento es evaluar si un sistema de generación de kernels GPU
basado en un Modelo de Lenguaje Pequeño (SLM) augmentado con decodificación restringida
por gramática (GBNF) y un ciclo de auto-reparación puede igualar o superar la tasa de
éxito de modelos frontier de mayor escala que utilizan \textit{prompting} directo sin
restricciones estructurales.

\subsection*{Preguntas de investigación}

\begin{itemize}
  \item \textbf{¿Qué problema se evalúa?} La generación automática de kernels Triton GPU
    correctos y compilables a partir de operaciones PyTorch, comparando un método
    especializado basado en SLM contra modelos baseline de mayor escala.

  \item \textbf{¿Qué se desea demostrar?} Que la combinación de decodificación restringida
    por gramática + ciclo de auto-reparación permite que un SLM alcance tasas de
    compilación y corrección comparables a modelos con órdenes de magnitud más parámetros.

  \item \textbf{¿Qué métodos se comparan?}
    \begin{itemize}
      \item \textbf{Método propuesto:} SLM (Llama-3.1-8B-Instruct via Fireworks AI) con
        pipeline especializado: gramática GBNF + validación en tres capas + auto-reparación.
      \item \textbf{Baseline 1:} Claude Opus 4.7 (Anthropic) — prompting directo, sin
        restricciones estructurales.
      \item \textbf{Baseline 2:} GPT-4o (OpenAI) — prompting directo, sin restricciones
        estructurales.
    \end{itemize}

  \item \textbf{Justificación de baselines:} Se seleccionaron dos modelos frontier
    ampliamente utilizados en la industria que representan el estado del arte en
    generación de código sin ajuste de dominio. La comparación es relevante porque
    evidencia el valor añadido por la infraestructura especializada frente a la capacidad
    bruta del modelo.
\end{itemize}

% -----------------------------------------------------------------------------
\section{Hipótesis}

Se definen hipótesis para las dos métricas principales: tasa de compilación y tasa de
corrección numérica.

\subsection{Hipótesis H1 — Tasa de Compilación}

\begin{itemize}
  \item $H_0^{(1)}$: No existe diferencia significativa en la tasa de compilación
    entre el método propuesto y Claude Opus 4.7.
    \[H_0^{(1)}: p_{\text{propuesto}} = p_{\text{Claude}}\]

  \item $H_1^{(1)}$: El método propuesto obtiene una tasa de compilación
    significativamente mayor o igual a la de Claude Opus 4.7.
    \[H_1^{(1)}: p_{\text{propuesto}} \geq p_{\text{Claude}}\]
\end{itemize}

\subsection{Hipótesis H2 — Tasa de Corrección Numérica}

\begin{itemize}
  \item $H_0^{(2)}$: No existe diferencia significativa en la tasa de corrección
    numérica entre el método propuesto y GPT-4o.
    \[H_0^{(2)}: p_{\text{propuesto}} = p_{\text{GPT-4o}}\]

  \item $H_1^{(2)}$: El método propuesto obtiene una tasa de corrección numérica
    significativamente mayor o igual a la de GPT-4o.
    \[H_1^{(2)}: p_{\text{propuesto}} \geq p_{\text{GPT-4o}}\]
\end{itemize}

\subsection{Hipótesis H3 — Intentos hasta Éxito}

\begin{itemize}
  \item $H_0^{(3)}$: No existe diferencia en el número promedio de intentos LLM
    hasta obtener un kernel válido entre el método propuesto y los baselines.

  \item $H_1^{(3)}$: El método propuesto requiere un número similar o menor de
    intentos que los baselines para producir un kernel válido, gracias a la
    restricción estructural de la gramática.
\end{itemize}

% -----------------------------------------------------------------------------
\section{Variables Experimentales}

\subsection{Variables Independientes}

Las variables independientes definen cómo está configurado cada método de generación:

\begin{table}[H]
\centering
\caption{Variables independientes por método}
\begin{tabular}{llll}
\toprule
\textbf{Variable} & \textbf{Método Propuesto} & \textbf{Baseline 1 (Claude)} & \textbf{Baseline 2 (GPT-4o)} \\
\midrule
Modelo base & Llama-3.1-8B & claude-opus-4-7 & gpt-4o \\
Parámetros (aprox.) & 8B & $\sim$200B+ & $\sim$200B+ \\
Proveedor & Fireworks AI & Anthropic & OpenAI \\
Temperatura & 0.0 & 0.0 & 0.0 \\
Gramática GBNF & Sí & No & No \\
Auto-reparación & Sí (máx. 3 intentos) & No & No \\
Prompt & Sistema + usuario & Solo usuario & Solo usuario \\
\bottomrule
\end{tabular}
\end{table}

\subsection{Variables Dependientes}

Las variables dependientes son las métricas que se miden para realizar la comparación:

\begin{itemize}
  \item \textbf{Tasa de compilación} ($CR$): Proporción de kernels que pasan la
    validación de compilación GPU.
  \item \textbf{Tasa de corrección numérica} ($NR$): Proporción de kernels que pasan
    \texttt{torch.allclose} contra baseline PyTorch.
  \item \textbf{Diferencia máxima absoluta} (\textit{max\_diff}): Error numérico
    máximo entre la salida del kernel y la referencia PyTorch.
  \item \textbf{Intentos hasta éxito} ($A$): Número de llamadas LLM requeridas para
    producir un kernel válido.
  \item \textbf{Speedup} ($S$): Cociente entre el tiempo de ejecución PyTorch y el
    tiempo del kernel generado.
    \[S = \frac{t_{\text{PyTorch}}}{t_{\text{kernel}}}\]
\end{itemize}

\subsection{Variables Controladas}

Las siguientes variables se mantienen fijas en todos los experimentos para garantizar
comparabilidad:

\begin{itemize}
  \item \textbf{Hardware:} NVIDIA T4 (Google Colab), misma sesión por método.
  \item \textbf{Benchmark:} Mismo conjunto de operaciones PyTorch para todos los métodos.
  \item \textbf{Temperatura:} $T = 0.0$ en todos los modelos (salidas deterministas).
  \item \textbf{Número de operaciones evaluadas:} $N = 30$ operaciones por método.
  \item \textbf{Prompt de sistema:} El prompt de sistema del método propuesto incluye
    reglas explícitas de Triton; los baselines reciben el mismo enunciado de tarea.
  \item \textbf{Semilla aleatoria de inputs:} \texttt{torch.manual\_seed(42)} para
    todos los tensores de prueba.
\end{itemize}

% -----------------------------------------------------------------------------
\section{Diseño Experimental}

\subsection{Estructura}

Se utiliza un diseño \textbf{within-subjects} (intra-sujeto), donde cada operación
PyTorch del benchmark es evaluada por los tres métodos. Esto permite comparaciones
directas por par de operación, reduciendo la variabilidad introducida por diferencias
en dificultad entre operaciones.

\begin{itemize}
  \item \textbf{Número de grupos:} 3 (método propuesto, Claude Opus 4.7, GPT-4o).
  \item \textbf{Número de operaciones por grupo:} 30.
  \item \textbf{Total de llamadas LLM:} $\geq 90$ (más intentos de reparación).
\end{itemize}

\subsection{Benchmark de Operaciones}

El benchmark incluye operaciones representativas de tres categorías de complejidad:

\begin{table}[H]
\centering
\caption{Categorías del benchmark}
\begin{tabular}{llp{7cm}}
\toprule
\textbf{Categoría} & \textbf{\# ops} & \textbf{Ejemplos} \\
\midrule
Elementwise simple & 10 & \texttt{out = x * y + z}, \texttt{out = relu(x)}, \texttt{out = x + y} \\
Reducción & 10 & \texttt{out = sum(x)}, \texttt{out = max(x, dim=0)} \\
Compuesta & 10 & \texttt{out = softmax(x)}, \texttt{out = layer\_norm(x, w, b)} \\
\bottomrule
\end{tabular}
\end{table}

\subsection{Justificación del Diseño}

El diseño within-subjects se eligió porque:
\begin{enumerate}
  \item Elimina la variabilidad entre operaciones al comparar los tres métodos sobre
    exactamente los mismos inputs.
  \item Requiere menor tamaño de muestra para el mismo poder estadístico.
  \item Permite comparaciones pareadas que tienen mayor sensibilidad que comparaciones
    de grupos independientes.
\end{enumerate}

% -----------------------------------------------------------------------------
\section{Métricas}

\begin{table}[H]
\centering
\caption{Definición de métricas}
\begin{tabular}{p{3cm}p{4cm}p{4cm}p{3cm}}
\toprule
\textbf{Métrica} & \textbf{Fórmula / Criterio} & \textbf{Interpretación} & \textbf{Limitación} \\
\midrule
Tasa de compilación ($CR$) &
  $CR = \frac{\text{\# kernels compilados}}{\text{\# kernels generados}}$ &
  Mayor es mejor. $CR=1.0$ indica que todos los kernels compilaron. &
  No distingue calidad de kernels que compilan pero son incorrectos. \\
\midrule
Tasa de corrección ($NR$) &
  $NR = \frac{\text{\# kernels correctos}}{\text{\# kernels compilados}}$ &
  Mayor es mejor. Requiere GPU para ejecutar. &
  Depende de tolerancias \texttt{rtol=1e-3, atol=1e-3}; puede no capturar errores numéricos sutiles. \\
\midrule
Max diff &
  $\max |y_{\text{kernel}} - y_{\text{pytorch}}|$ &
  Menor es mejor. $0.0$ indica exactitud perfecta. &
  Valor único; puede ocultar distribución del error. \\
\midrule
Intentos ($A$) &
  Número de llamadas LLM hasta pasar todas las capas de validación. &
  Menor es mejor. $A=1$ indica éxito en primer intento. &
  Refleja tanto calidad del modelo como severidad del benchmark. \\
\midrule
Speedup ($S$) &
  $S = t_{\text{pytorch}} / t_{\text{kernel}}$ &
  $S > 1$: kernel más rápido que PyTorch. $S < 1$: overhead. &
  Varía por GPU, tamaño de tensor, calentamiento. \\
\bottomrule
\end{tabular}
\end{table}

% -----------------------------------------------------------------------------
\section{Tamaño del Efecto Esperado}

Se espera un efecto mediano ($d = 0.5$, Cohen) en la diferencia de tasas de compilación,
basado en la suposición de que la gramática GBNF elimina sistemáticamente errores de
estructura (imports faltantes, decorador \texttt{@triton.jit} ausente) que los baselines
cometen en una fracción de los casos.

Para tasas de corrección numérica, se espera un efecto menor ($d \approx 0.3$) ya que
este atributo depende más de la calidad de razonamiento del modelo que de las restricciones
estructurales.

% -----------------------------------------------------------------------------
\section{Análisis de Poder}

\begin{itemize}
  \item \textbf{Nivel de significancia:} $\alpha = 0.05$
  \item \textbf{Poder estadístico objetivo:} $1 - \beta = 0.80$
  \item \textbf{Efecto esperado:} $d = 0.5$ (mediano, Cohen)
  \item \textbf{Tamaño de muestra estimado:} Para una prueba de McNemar con $\alpha=0.05$
    y poder $0.8$, se requieren aproximadamente $N = 28$ pares. Se utilizarán 30 para
    proveer un margen de seguridad.
\end{itemize}

\[
n \approx \frac{(z_{\alpha/2} + z_{\beta})^2}{(\text{efecto esperado})^2}
= \frac{(1.96 + 0.84)^2}{0.5^2} \approx 31.4
\]

Por lo tanto, $N = 30$ operaciones por método es suficiente dado el efecto esperado.

% -----------------------------------------------------------------------------
\section{Plan Estadístico}

\subsection{Comparaciones Planeadas}

\begin{enumerate}
  \item Método propuesto vs.\ Claude Opus 4.7 en tasa de compilación.
  \item Método propuesto vs.\ GPT-4o en tasa de compilación.
  \item Método propuesto vs.\ Claude Opus 4.7 en tasa de corrección numérica.
  \item Método propuesto vs.\ GPT-4o en tasa de corrección numérica.
  \item Método propuesto vs.\ ambos baselines en número de intentos hasta éxito.
  \item Método propuesto vs.\ ambos baselines en speedup (sólo kernels correctos).
\end{enumerate}

\subsection{Pruebas Estadísticas}

\begin{table}[H]
\centering
\caption{Pruebas estadísticas por métrica}
\begin{tabular}{lll}
\toprule
\textbf{Métrica} & \textbf{Prueba} & \textbf{Justificación} \\
\midrule
Tasa de compilación / corrección & Prueba de McNemar &
  Datos binarios pareados (mismo benchmark) \\
Intentos hasta éxito & Prueba de Wilcoxon signed-rank &
  Datos ordinales, no se asume normalidad \\
Speedup & Prueba de Wilcoxon signed-rank &
  Distribución de speedup generalmente no normal \\
\bottomrule
\end{tabular}
\end{table}

\subsection{Regla de Decisión}

Se rechazará $H_0$ si $p < \alpha = 0.05$ después de aplicar la corrección de
Bonferroni para comparaciones múltiples.

Con 6 comparaciones planeadas (sección 8.1), el umbral ajustado es:
\[\alpha_{\text{ajustado}} = \frac{0.05}{6} \approx 0.0083\]

\subsection{Intervalos de Confianza}

Se reportarán intervalos de confianza al 95\% para:
\begin{itemize}
  \item Diferencia de proporciones en tasa de compilación y corrección.
  \item Mediana del speedup por método.
  \item Mediana de intentos hasta éxito por método.
\end{itemize}

\subsection{Corrección por Comparaciones Múltiples}

Se aplicará corrección de Bonferroni dado que se realizan 6 comparaciones planeadas.
Si se detectaran comparaciones no planeadas post-hoc, se documentarán explícitamente
y se considerarán exploratorias.

% -----------------------------------------------------------------------------
\section{Reproducibilidad Experimental}

\subsection{Hardware}

\begin{table}[H]
\centering
\caption{Entorno de hardware}
\begin{tabular}{ll}
\toprule
\textbf{Componente} & \textbf{Especificación} \\
\midrule
GPU & NVIDIA T4 (Google Colab) \\
VRAM & 15 GB \\
CPU & Intel Xeon (Colab) \\
RAM & 12.7 GB \\
Sistema Operativo & Ubuntu 22.04 (Colab runtime) \\
\bottomrule
\end{tabular}
\end{table}

\subsection{Entorno de Software}

\begin{table}[H]
\centering
\caption{Versiones de software}
\begin{tabular}{ll}
\toprule
\textbf{Librería} & \textbf{Versión} \\
\midrule
Python & 3.11.x \\
PyTorch & 2.x \\
Triton & 3.x \\
fireworks-ai & $\geq$ 0.15.0 \\
python-dotenv & latest \\
\bottomrule
\end{tabular}
\end{table}

\textit{Nota: Las versiones exactas se reportarán en la Parte II después de la ejecución.}

\subsection{Modelos Utilizados}

\begin{table}[H]
\centering
\caption{Modelos evaluados}
\begin{tabular}{lllll}
\toprule
\textbf{Rol} & \textbf{Modelo} & \textbf{Versión/ID} & \textbf{Proveedor} & \textbf{Parámetros} \\
\midrule
Método propuesto & Llama 3.1 & \texttt{llama-v3p1-8b-instruct} & Fireworks AI & 8B \\
Baseline 1 & Claude Opus 4.7 & \texttt{claude-opus-4-7} & Anthropic & $\sim$200B+ \\
Baseline 2 & GPT-4o & \texttt{gpt-4o} & OpenAI & $\sim$200B+ \\
\bottomrule
\end{tabular}
\end{table}

\subsection{Configuración del Método Propuesto}

\begin{itemize}
  \item \textbf{Temperatura:} 0.0
  \item \textbf{Gramática GBNF:} Activa — fuerza imports + \texttt{@triton.jit} a nivel de token.
  \item \textbf{Max tokens:} 4096
  \item \textbf{Intentos máximos:} 3
  \item \textbf{Prompt de sistema:} Fijo (mismo para todas las operaciones).
\end{itemize}

\subsection{Repositorio}

\begin{itemize}
  \item \textbf{GitHub:} \url{https://github.com/casrhub/llm-gpu-kernel-generation}
  \item \textbf{Commit de referencia:} \texttt{ae1777a} (o el más reciente al momento
    de ejecutar experimentos)
  \item \textbf{Archivo de configuración:} \texttt{experiment\_config.json}
    (por crear antes de ejecutar)
\end{itemize}

% -----------------------------------------------------------------------------
\section{Amenazas a la Validez (Pre-registro)}

\subsection{Validez Interna}

\begin{itemize}
  \item \textbf{Diferencias de hardware entre corridas:} Colab asigna GPU dinámicamente.
    Se mitigará ejecutando todos los métodos en la misma sesión de Colab cuando sea
    posible, o documentando el tipo de GPU en cada corrida.

  \item \textbf{No-determinismo residual:} Aunque $T=0.0$, algunos modelos pueden
    presentar variabilidad por precisión de punto flotante en inference distribuida.
    Se mitigará ejecutando cada operación 3 veces y reportando la moda.

  \item \textbf{Selección del benchmark:} Si el benchmark favorece operaciones simples
    donde la gramática ayuda más, el método propuesto podría estar sobreestimado.
\end{itemize}

\subsection{Validez Externa}

\begin{itemize}
  \item \textbf{Generalización limitada:} El benchmark incluye 30 operaciones.
    Las conclusiones pueden no generalizarse a operaciones más complejas
    (transformers completos, operaciones matriciales avanzadas).

  \item \textbf{Hardware específico:} Los resultados de speedup son específicos a
    NVIDIA T4. En otras GPUs (A100, H100) la relación podría cambiar.
\end{itemize}

\subsection{Validez de Constructo}

\begin{itemize}
  \item \textbf{Tasa de compilación como proxy de calidad:} Un kernel puede compilar
    pero ser numéricamente incorrecto o ineficiente. La compilación no implica
    utilidad real.

  \item \textbf{Diferencia de escala entre modelos:} Los baselines tienen órdenes de
    magnitud más parámetros. El speedup en tasa de compilación puede explicarse
    únicamente por la restricción estructural (gramática), no por la calidad del SLM.
\end{itemize}

% =============================================================================
% PARTE II — RESULTADOS Y ANÁLISIS
% =============================================================================
\newpage
\part{Resultados y Análisis}
\textit{Esta sección se completará después de ejecutar los experimentos comparativos.}

% -----------------------------------------------------------------------------
\section{Resultados Descriptivos}

% TODO: Completar con datos reales después de ejecutar experimentos

\begin{table}[H]
\centering
\caption{Resumen de resultados por método}
\begin{tabular}{lcccccc}
\toprule
\textbf{Método} & \textbf{N} & \textbf{CR (\%)} & \textbf{NR (\%)} &
  $\bar{A}$ & $\sigma_A$ & \textbf{Speedup $\bar{S}$} \\
\midrule
SLM + GBNF + reparación & 30 & -- & -- & -- & -- & -- \\
Claude Opus 4.7 & 30 & -- & -- & -- & -- & -- \\
GPT-4o & 30 & -- & -- & -- & -- & -- \\
\bottomrule
\end{tabular}
\end{table}

\textit{CR = tasa de compilación, NR = tasa de corrección numérica, $\bar{A}$ =
media de intentos, $\sigma_A$ = desviación estándar de intentos.}

% -----------------------------------------------------------------------------
\section{Visualización de Datos}

% TODO: Insertar figuras después de ejecutar experimentos

\begin{figure}[H]
\centering
% \includegraphics[width=0.7\textwidth]{figures/compilation_rate_bar.pdf}
\framebox[0.7\textwidth]{\rule{0pt}{4cm} [Gráfica de barras — Tasa de compilación por método]}
\caption{Tasa de compilación por método. Barras de error representan
  intervalos de confianza al 95\%.}
\end{figure}

\begin{figure}[H]
\centering
% \includegraphics[width=0.7\textwidth]{figures/speedup_boxplot.pdf}
\framebox[0.7\textwidth]{\rule{0pt}{4cm} [Box plot — Speedup por método]}
\caption{Distribución de speedup por método. Solo se incluyen kernels que
  pasaron la validación de corrección numérica. La línea punteada en $S=1$
  indica paridad con PyTorch.}
\end{figure}

% -----------------------------------------------------------------------------
\section{Inferencia Estadística}

% TODO: Completar con p-values y estadísticos reales

\begin{table}[H]
\centering
\caption{Resultados de pruebas estadísticas}
\begin{tabular}{lllllll}
\toprule
\textbf{Comparación} & \textbf{Métrica} & \textbf{Prueba} &
  \textbf{Estadístico} & \textbf{p-value} & \textbf{p ajustado} & \textbf{d Cohen} \\
\midrule
Propuesto vs.\ Claude & CR & McNemar & -- & -- & -- & -- \\
Propuesto vs.\ GPT-4o & CR & McNemar & -- & -- & -- & -- \\
Propuesto vs.\ Claude & NR & McNemar & -- & -- & -- & -- \\
Propuesto vs.\ GPT-4o & NR & McNemar & -- & -- & -- & -- \\
Propuesto vs.\ Claude & Intentos & Wilcoxon & -- & -- & -- & -- \\
Propuesto vs.\ GPT-4o & Intentos & Wilcoxon & -- & -- & -- & -- \\
\bottomrule
\end{tabular}
\end{table}

\subsection*{Interpretación (plantilla)}

\begin{itemize}
  \item \textbf{¿Se rechazó $H_0$?} [Completar]
  \item \textbf{Evidencia estadística:} [Completar con p-values]
  \item \textbf{Tamaño del efecto:} [Pequeño / Mediano / Grande según d de Cohen]
  \item \textbf{Consistencia entre operaciones:} [Completar]
\end{itemize}

% -----------------------------------------------------------------------------
\section{Comparación Crítica}

% TODO: Completar con análisis real

\textit{[Discutir si los resultados fueron consistentes entre las tres categorías del
benchmark (elementwise, reducción, compuesta). Analizar si el método propuesto tuvo
ventaja particular en alguna categoría. Considerar si las diferencias observadas pueden
explicarse parcialmente por la diferencia de escala entre modelos.]}

% -----------------------------------------------------------------------------
\section{Amenazas Observadas a la Validez}

\begin{table}[H]
\centering
\caption{Amenazas detectadas durante la ejecución}
\begin{tabular}{p{3cm}p{5cm}lp{4cm}}
\toprule
\textbf{Amenaza} & \textbf{Descripción} & \textbf{Clasificación} & \textbf{Mitigación aplicada} \\
\midrule
Asignación dinámica de GPU en Colab &
  Distintas sesiones pueden recibir diferentes modelos de GPU, afectando speedup. &
  Interna &
  Documentar tipo de GPU por sesión. \\
Validador con bugs &
  Errores en el pipeline de validación (exec() vs.\ temp file) generaron falsos negativos
  durante el desarrollo, afectando la interpretación inicial. &
  Interna &
  Correcciones aplicadas y versionadas en git. \\
Benchmark pequeño &
  30 operaciones pueden ser insuficientes para detectar diferencias en operaciones
  complejas (compuesta). &
  Externa &
  Análisis separado por categoría. \\
\midrule
\multicolumn{4}{c}{\textit{[Agregar amenazas adicionales detectadas durante experimentos]}} \\
\bottomrule
\end{tabular}
\end{table}

% -----------------------------------------------------------------------------
\section{Discusión}

% TODO: Completar después de ejecutar experimentos

\textit{[Guía: discutir si la gramática GBNF realmente aportó valor o si los baselines
también habrían compilado al primer intento. Analizar si el ciclo de auto-reparación
fue necesario. Evaluar si la diferencia de resultados entre el SLM y los modelos frontier
se explica por el tamaño del modelo o por la infraestructura especializada. Discutir
qué limitaciones quedaron abiertas.]}

% -----------------------------------------------------------------------------
\section{Conclusión}

% TODO: Completar después de ejecutar experimentos

\textit{[Guía: responder explícitamente cada pregunta de la rúbrica: ¿se rechazó $H_0$?
¿Hubo evidencia estadística suficiente? ¿La diferencia fue consistente? ¿El tamaño del
efecto fue relevante? ¿La mejora justificaría usar el método en producción? ¿Qué
limitaciones identificaron? ¿Podría otro equipo reproducir el experimento?]}

% =============================================================================
\end{document}
